\documentclass[11pt]{article}
\usepackage[margin=1in]{geometry}
\usepackage{amsmath,amssymb}
\usepackage{booktabs}
\usepackage{array}
\usepackage{hyperref}
\usepackage{enumitem}
\usepackage{xcolor}
\hypersetup{colorlinks=true,linkcolor=blue,citecolor=blue,urlcolor=blue}

\title{\textbf{Construction of the Event Rainfall and Antecedent Wetness Variables} \\[4pt]
\large Methods Note --- from mechanism to a computable, calibrated metric}
\author{MMSD Wetland--Flood Effectiveness Study (rainfall\_v2)}
\date{July 2026}

\begin{document}
\maketitle

\begin{abstract}
This note documents, end to end, how we construct the event--level rainfall predictors used in
the event--based flood--response models --- basin--mean event rainfall $P_e$, event rainfall
volume $V_e$, storm intensity, and antecedent basin wetness --- and why each modeling choice is
made, with the antecedent-wetness construction anchored to the U.S.\ Weather Bureau's classic
antecedent precipitation index (API) methodology (Kohler \& Linsley, 1951). Starting from a
radar-based basin-average rainfall series, we compute event-window rainfall depth and volume, an
event runoff coefficient, storm intensity, and a decay-weighted cumulative rainfall term. We
identify and correct an off-by-one indexing error in the original antecedent-wetness lookup, and
adopt an ``effective rainfall'' formulation that includes the event's own start-day rainfall in
the decay-weighted sum, consistent with the Weather Bureau's logarithmic-recession antecedent
index but extended to the day of interest. We then describe how these event-level rainfall
variables are joined with gauge-level wetland-effectiveness and control variables to build the
modeling panel, and how they enter as predictors and controls in the phenomenological models fit
downstream.
\end{abstract}

\section{Motivation: turning a rainfall record into event-level predictors}

The scientific question addressed by the broader project is whether upstream wetland
configuration is associated with the flood response measured at a gauge. Answering it requires
controlling for the rainfall that produced each flood event, since runoff response is strongly
driven by how much rain fell, how intense it was, and how wet the basin already was before the
storm began. This note documents the construction of that rainfall side of the model: for every
gauge $j$ and storm event $e$ (as delineated upstream in the event-identification step), we
build

\begin{itemize}[nosep]
  \item $P_{e}$ --- basin-mean rainfall depth over the event window;
  \item $V_{e}$ --- the corresponding rainfall volume, used to form an event runoff coefficient;
  \item $\text{max\_daily}$ --- peak one-day rainfall intensity within the event window;
  \item $P_{\text{eff}}$ --- a decay-weighted cumulative rainfall term representing antecedent
        basin wetness at the start of the event.
\end{itemize}

The guiding principle, as in the companion wetland-effectiveness note, is to avoid inventing
variables: the antecedent-wetness construction is borrowed from an authoritative published
method rather than an arbitrary index.

\section{Literature review: the antecedent precipitation index}

The construction of $P_{\text{eff}}$ follows the U.S.\ Weather Bureau's Research Paper No.\ 34,
\emph{Predicting the Runoff from Storm Rainfall} (Kohler \& Linsley, 1951), which documents the
technique used at the Weather Bureau's River Forecast Centers for evaluating antecedent
conditions in flood forecasting. The paper reviews several candidate moisture indices --- days
since last rain, discharge at the beginning of the storm, and antecedent precipitation --- and
concludes that antecedent precipitation is the most broadly applicable, ``provided the index is
properly derived and is used in conjunction with season of the year or temperature''
(Kohler \& Linsley, 1951, p.\ 2).

The paper defines the antecedent precipitation index (API) generally as
\begin{equation}
  I = b_1 P_1 + b_2 P_2 + b_3 P_3 + \dots + b_i P_i,
  \label{eq:api-general}
\end{equation}
where $P_i$ is the precipitation occurring $i$ days prior to the storm under consideration and
$b_i$ is a weight assumed to decay with time. The paper notes an important refinement: rather
than a reciprocal weight ($b_i = 1/i$), a \emph{logarithmic recession} is preferable for
day-to-day computation, i.e.\ during a dry spell
\begin{equation}
  I_t = I_0\, k^{t}, \qquad k^{}\text{(unitless recession factor)},
  \label{eq:api-recession}
\end{equation}
so that each day's index is simply the previous day's index multiplied by $k$, with that day's
rainfall added when it occurs. The paper reports that $k$ is not a critical parameter and that
values in the range $0.85$--$0.90$ are appropriate across most of the eastern and central United
States. Our implementation adopts $k = 0.90$, at the upper (slower-decay) end of that range, and
a 30-day summation window.

\section{Design decision: antecedent index vs.\ effective cumulative rainfall}

Equation~\eqref{eq:api-general}, exactly as stated in Kohler \& Linsley (1951), sums
precipitation from $i=1$ days prior onward --- it deliberately excludes the storm's own rainfall,
since the index is meant to represent basin wetness \emph{before} the event begins. Our first
implementation followed this definition exactly.

In reviewing the variable for use as a model control, we adopted a related but distinct
convention that is also long-standing in hydrologic practice: an \emph{effective cumulative
rainfall} term that folds the event's own start-day rainfall into the same decay-weighted sum,
i.e.\ extending the summation to include an $i=0$ term at full weight:
\begin{equation}
  P_{\text{eff}}(t) = P(t) + k\,P(t-1) + k^2 P(t-2) + \dots + k^{30} P(t-30).
  \label{eq:peff}
\end{equation}
This is a deliberate extension beyond the literal Kohler \& Linsley definition (which is
strictly $i \geq 1$): we flag this explicitly rather than attribute it to the cited paper, since
the paper's own index is antecedent-only. The practical effect of the change is that the
variable now represents cumulative decay-weighted rainfall \emph{through and including} the
event's start day, rather than strictly the wetness state existing immediately before it. We
made this choice because it is the more commonly used convention in comparable modeling
pipelines and because, in combination with the correction described in \S4, it removes an
ambiguity in what date the lookup should reference.

\section{Correcting an indexing error in the antecedent-wetness lookup}

While reviewing the event-pairing code, we identified an off-by-one error in how the
antecedent-wetness value was retrieved for each event. The index series itself
(Eq.~\eqref{eq:api-general}, before the change described in \S3) was computed correctly, but the
per-event lookup subtracted an additional day before reading the series:
\begin{quote}
\ttfamily
prev = d0 - 1 day \\
api\_v = API series value at prev
\end{quote}
Because the original API series already excludes the lookup day's own rainfall by construction,
subtracting a further day meant the single most heavily weighted, most recent term (the day
immediately preceding the event) was silently dropped from every event's antecedent-wetness
value. In effect, the stored variable reflected basin wetness as of two days before the event
rather than one.

Adopting the effective-rainfall formulation of Eq.~\eqref{eq:peff} resolves this cleanly: the
lookup is now taken directly at the event's start date $d_0$, with no additional day subtracted,
since $P_{\text{eff}}(d_0)$ is now defined to include $d_0$'s own rainfall by design. The fix
therefore consists of (i) extending the summation in the index series to include the $i=0$ term,
and (ii) removing the extra day offset in the event-level lookup. No other rainfall variable
($P_e$, $V_e$, runoff coefficient, or intensity) was affected by this error, since each of those
is computed by direct summation over the event window rather than through the lagged index
series.

\section{Constructing each rainfall variable}

\subsection{Basin-mean daily rainfall}
Daily rainfall is drawn from NCEP Stage IV radar-gauge precipitation (IEM IEMRE web service,
0.125\textdegree\ grid, 1997--present), area-averaged over each catchment by sampling a dense
grid of interior points, snapping each to its nearest 0.125\textdegree\ lattice node, and
weighting nodes by their share of interior sample points. This produces one basin-mean daily
rainfall series per gauge.

\subsection{Event rainfall depth and volume ($P_e$, $V_e$)}
For each event window $[t_{\text{start}}, t_{\text{end}}]$ (as delineated in the upstream
event-identification step), $P_e$ is the sum of basin-mean daily rainfall over the window, and a
separate sum to the event's peak, $P_{\text{topeak}}$, is retained. Rainfall volume is
$V_e = (P_e / 1000)\times A$, where $A$ is catchment area in m$^2$. Dividing the event's measured
quick-flow runoff volume by $V_e$ gives an event runoff coefficient.

\subsection{Storm intensity}
Peak one-day rainfall intensity within the event window is retained (\texttt{max\_daily\_mm}) as
a simple intensity proxy, since two events of equal depth and equal antecedent wetness can
produce different peak responses depending on how concentrated the rainfall was in time.

\subsection{Antecedent / effective cumulative rainfall ($P_{\text{eff}}$)}
Computed per Eq.~\eqref{eq:peff} with $k=0.90$ over a 30-day window, looked up directly at each
event's start date as described in \S4. This is the term intended to control for basin wetness
entering each storm.

\section{Downstream use: assembling the modeling panel and fitting the models}

The event-level rainfall variables documented here are not consumed directly by the final
models; they first pass through two further assembly and modeling steps, summarized here for
completeness.

\subsection{Panel assembly}
A separate script joins the per-event rainfall table (one row per gauge-event, containing $P_e$,
$V_e$, runoff coefficient, $P_{\text{eff}}$, and intensity) with two gauge-level tables: a
wetland-effectiveness table (contributing $W$-variants and wetland storage $S_w$; see the
companion wetland-effectiveness note) and a static-controls table (drainage area, slopes,
saturated conductivity, soil sand fraction, available water storage, and imperviousness
interpolated to each event's year). It then derives modeling terms directly from the rainfall
and wetland variables already constructed --- most notably a saturation ratio $V_e/S_w$ relating
event rainfall volume to gauge-level wetland storage capacity, a winter-season flag, a flag for
runoff coefficients exceeding one (a data-quality check), and log-transformed versions of peak
discharge, $P_e$, $V_e$, drainage area, wetland storage, and wetland effectiveness. Each
gauge-event is assigned a unique identifier, and a usability flag restricts the panel to events
with $P_e > 1$~mm and a finite runoff coefficient. The result is a single gauge-event panel used
as the modeling dataset, alongside a gauge-level collinearity diagnostic among the candidate
predictors.

\subsection{Model fitting}
The panel is used to fit a sequence of four nested phenomenological models relating a flood
response outcome (log peak discharge, Richards--Baker flashiness, or hydrograph width) to
standardized rainfall ($P$, built from $\log P_e$), wetland effectiveness ($W$, an
intensity-weighted wetland fraction), their interaction ($W{:}P$, a peak-shaving term), a
saturation term ($S = \log(V_e/S_w)$), and its interaction with wetland effectiveness
($W{:}S$), together with gauge-level controls. The antecedent-wetness variable constructed here
enters as one of those controls (alongside drainage area, imperviousness, saturated conductivity,
channel slope, and a winter flag), so that the wetland and rainfall-interaction terms of primary
interest are estimated net of antecedent basin wetness. Models are estimated both as
gauge-random-intercept mixed models and as OLS with gauge-clustered standard errors, and are
subjected to two falsification checks: a gauge-level wetland-effectiveness permutation placebo,
and leave-one-gauge-out refitting to assess coefficient stability.

\section{What remains}

Two items are worth flagging for completeness, mirroring the candor of the companion
wetland-effectiveness note. First, the switch from a strict antecedent-only index to an
effective-cumulative-rainfall formulation (\S3) is a deliberate deviation from the literal
Kohler \& Linsley (1951) definition, adopted for convenience and internal consistency rather than
because the cited paper itself recommends including the event's own start-day rainfall; readers
should not treat Eq.~\eqref{eq:peff} as a verbatim reproduction of the Weather Bureau method.
Second, the recession factor $k=0.90$ and the 30-day window are within, but at the upper edge of,
the range the source literature reports as broadly applicable ($k=0.85$--$0.90$); we have not
independently recalibrated $k$ against the observed events used in this project, and doing so
against event-level runoff response would be a natural next check.

\section*{References}
\begin{enumerate}[label={[\arabic*]},leftmargin=2.2em]
  \item Kohler, M.\ A., \& Linsley, R.\ K.\ (1951). \emph{Predicting the Runoff from Storm
  Rainfall}. U.S.\ Weather Bureau Research Paper No.\ 34. Washington, D.C.: U.S.\ Department of
  Commerce, Weather Bureau.
  \item MMSD Wetland--Flood study (2026). \emph{Construction of the cumulative wetland
  effectiveness variable $W$: methods and literature review} (companion note).
  \item Sherman, L.\ K.\ (1932). Streamflow from rainfall by the unit-graph method.
  \emph{Engineering News-Record}, 108(14), 501--505.
  \item Linsley, R.\ K., Kohler, M.\ A., \& Paulhus, J.\ L.\ H.\ (1949). \emph{Applied Hydrology}.
  New York: McGraw-Hill.
\end{enumerate}

\end{document}
